The following concept is crucial for several results.

\begin{defn}\label{thm:defqinc} A sequence $(b_n)$ of positive numbers is called eventually $q$-increasing, if there exist a number $0<q<1$ and $n_0\in\N$ such that
\begin{gather}\label{eq:qinc}
\frac{b_{n-1}}{b_n}\le q,\qquad n\ge n_0.
\end{gather}
For $n_0=1$ we drop ``eventually''.
\end{defn}

Note that if $(b_n)$ is unbounded and eventually log-convex, then it is eventually $q$-increasing with $q=b_{n_0-1}/b_{n_0}<1$, where $n_0$ is so large that $(b_n)$ is log-convex and strictly increasing for $n \ge n_0$.

We recall that an indeterminate moment problem has an order $\rho$ satisfying $0 \le\rho\le 1$. This is the common order of the four functions of the Nevanlinna matrix, see~\cite{B:P}.

\begin{prop}\label{thm:qinc-indet}\quad
\begin{enumerate}\itemsep=0pt
\item[$(i)$] The symmetric moment problem for an eventually $q$-increasing sequence $(b_n)$ is indeterminate and of order~$0$.

\item[$(ii)$] The symmetric moment problem for a positive sequence $(b_n)$ satisfying
\begin{gather}\label{eq:betan}
\frac{b_{n-1}}{b_n}\le {\rm e}^{-\beta/n},\qquad n\ge n_0, \qquad \beta>1,
\end{gather}
is indeterminate of order $\le 1/\beta$.
\end{enumerate}
\end{prop}

\begin{proof} Let $(Q_n)$ denote the sequence of polynomials of the second kind. By \cite[Theorem~1.2]{B:S2}, it suffices to prove that $\big(P_n^2(0)\big), \big(Q_n^2(0)\big)\in\ell^\alpha$ for all $0<\alpha\le 1$ in the first case and for all $\alpha>1/\beta$ in the second case.

 By the symmetry assumption we know that $P_{2n+1}(0)=Q_{2n}(0)=0$ and by \cite[Remark 4.5]{B:S2} we have
\begin{gather}\label{eq:P2n(0)}
P_{2n}(0)=(-1)^n\frac{b_0b_2\cdots b_{2n-2}}{b_1b_3\cdots b_{2n-1}},\qquad Q_{2n+1}(0)=(-1)^n\frac{b_1b_3\cdots b_{2n-1}}{b_0b_2\cdots b_{2n}}.
\end{gather}
For $n>n_0$ such that \eqref{eq:qinc} holds, we have
\begin{gather*}
P_{2n}^2(0)\le \left(\frac{b_0b_2\cdots b_{2n_0-2}}{b_1b_3\cdots b_{2n_0-1}}\right)^2q^{2(n-n_0)},\qquad
Q_{2n+1}^2(0)\le \left(\frac{b_1b_3\cdots b_{2n_0-1}}{b_0b_2\cdots b_{2n_0}}\right)^2 q^{2(n-n_0)},
\end{gather*}
and the first assertion follows.

For $n>n_0$ such that \eqref{eq:betan} holds, we have
\begin{gather*}
P_{2n}^2(0)\le \left(\frac{b_0b_2\cdots b_{2n_0-2}}{b_1b_3\cdots b_{2n_0-1}}\right)^2\exp\left(-2\beta\sum_{j=n_0}^{n-1}\frac{1}{2j+1}\right),
\end{gather*}
and a similar inequality for $Q_{2n+1}^2(0)$. Using the inequality
\begin{gather*}
\sum_{j=n_0}^{n-1}\frac{1}{2j+1}\ge \int_{n_0}^n\frac{{\rm d}x}{2x+1}=\frac12\log\frac{2n+1}{2n_0+1},
\end{gather*}
we get
\begin{gather*}
P_{2n}^2(0)\le \frac{C}{(2n+1)^\beta}
\end{gather*}
for a suitable constant $C$, and the second assertion follows.
\end{proof}
